\section{Introduction}

Researchers who study large collections of text work with samples, because the collections contain more text than people or models can process. An annotation team labels as many posts as its budget pays for. A large language model takes in only as much text as fits in its context window. Qualitative researchers stop interviewing at saturation, when new interviews stop adding new themes, which the literature places at about 9 to 17 interviews \citep{guest2006,hennink2022}. In each case a budget or a convention sets the size of the sample. If we read 15 posts from a corpus of 10{,}000, can we be sure we basically understand the corpus? Would 100 or 1{,}000 posts be enough?

Survey research answers the same question for an estimate about a population. Before the survey, a sample size formula tells the designer how many respondents a target margin of error requires \citep{cochran1977,kish1965}. After the survey, the margin of error tells the reader how far the estimate may lie from the population value. A sample of a corpus comes with neither. The question a reader of such a sample has in mind is how much of the corpus has a similar document among the documents read. This paper builds an estimate of that percentage and a sample size formula for it.

A text encoder turns each document into an embedding, a list of numbers, so that documents with similar meaning have embeddings close together. We define two documents as the same content when the cosine distance between their embeddings is smaller than a resolution $r$ chosen by the researcher. A researcher who needs the same specific claim chooses a small $r$, and one who needs only the same topic chooses a larger one. To give $r$ a meaning in content, two large language models judged pairs of documents at a range of distances.

With $r$ fixed, a corpus document is missed when no sampled document lies within $r$ of it. The missing mass is the percentage of the corpus that is missed. A researcher who holds only the sample cannot count the missed documents directly. Good and Turing faced the same problem with words, and estimated the probability that the next word of a text is one not yet seen from the words seen exactly once \citep{good1953}. \citet{maurer2022missing} extended this estimator from words to distances. Its input is what we call the isolation rate $G$, the percentage of sampled documents with no other sampled document within $r$. A sample of $n$ documents from a corpus of $N$ estimates the missing mass as $G\,(1 - n/N)$, where $1 - n/N$ is the finite population correction of survey sampling. Maurer develops the estimator in theory, and in a search of the literature in 2026 we found no work that applies it to text.

How fast the missing mass falls as the sample grows follows from a classical geometric result about nearest neighbors among randomly scattered points \citep{penrose2003rgg,levina2005}. The intrinsic dimension $d$ of a corpus is the number of independent ways in which its documents differ in meaning. We find that, on most of the corpora we tested, any two samples with the same product $n r^{d/2}$ leave the same percentage of their documents isolated. A larger sample and a looser resolution therefore trade against each other at a fixed rate set by $d$. Under this law, halving $r$ requires $2^{d/2}$ times as many documents, which for $d = 26$ is about 8{,}000 times as many. The dimension $d$ does for a corpus what the variance does for a survey population. A population with a larger variance needs more respondents for the same margin of error, and a corpus with a larger $d$ needs more documents for the same resolution.

We test the estimate and the formula on synthetic data with a known answer and on ten real corpora: three Reddit forums, five slices of 2020 tweets about the United States Census, and random samples of English posts from Truth Social and Bluesky in 2024. The paper makes four contributions.
\begin{enumerate}
\item A measure of what a sample of text has missed, the missing mass at a resolution chosen by the researcher, which the sample alone estimates through the isolation rate and the finite population correction of survey sampling.
\item A sample size formula, obtained by solving the law for $n$, with a pilot test that tells whether the formula can be trusted on a given corpus.
\item A way to give the resolution $r$ a meaning in content, by having large language models judge pairs of documents.
\item A measurement of how much of a corpus a sample chosen by $k$-means covers against a random sample of the same size. The chosen sample covers far more of the corpus on corpora that repeat themselves.
\end{enumerate}
